The experimental branches gathered in this part are not presented as privileged embodiments of Harmonic Information Theory. They are the branches that were already taking shape, for heterogeneous practical and research reasons, within the broader work from which HIT gradually became thinkable. Phideus and Beacon did not arise as applications of a finished doctrine. They are among the sites where the need for such a doctrine became increasingly difficult to ignore.

That contingency should not be mistaken for exclusivity. Many of the works gathered earlier as convergences across fields may also be read, from the perspective opened by HIT, as experimental branches that touch the same object without yet naming it in those terms. If researchers in other domains were to work explicitly through the hermeneutic proposed here, the number of possible experimental branches would have no obvious closure. The probes gathered in this part should therefore be read for what they are: not a privileged map of HIT's experimental horizon, and not anything like an exhaustive survey of it, but a first cluster of situated exploratory fronts through which the theory begins to undergo explicit experimental testing.
